\section{Le \textit{Livre-Journal} : description matérielle et histoire du document}

\subsection{Un petit carnet relié en tête}

Le \textit{Livre-Journal} est un carnet de poche relié en tête, comme un bloc-notes : on écrit le recto, puis on rabat le feuillet par-dessus la reliure pour écrire le verso. Le plat supérieur de la couverture, en papier rouge gaufré à losanges, porte une étiquette bleue où Hippolyte Combes a écrit le titre et les dates du carnet : \textit{Livre~-~Journal | Du 25~juillet 1916 | Au 30~avril 1917}. Les dimensions de l'original n'ont pas été relevées ; les pages, en hauteur, ont un rapport largeur/hauteur d'environ 0,68.

\begin{figure}[ht]
\centering
\begin{tikzpicture}[x=0.72mm,y=0.72mm,line width=0.4pt,font=\footnotesize]
  % --- carnet relié en tête ---
  \begin{scope}
    \fill[gray!8] (0,0) rectangle (34,48);
    \draw (0,0) rectangle (34,48);
    \fill[gray!45] (0,46.2) rectangle (34,48.8);
    \foreach \y in {41,37,33,29,25,21,17} {\draw[gray!70] (4,\y) -- (30,\y);}
    \node[font=\scriptsize\itshape] at (17,44) {recto};
    \node[align=center,font=\scriptsize] at (52,30) {le feuillet\\se rabat\\par-dessus\\la reliure};
    % feuillet rabattu (verso), au-dessus de la reliure
    \fill[gray!4] (0,48.8) rectangle (34,94);
    \draw (0,48.8) rectangle (34,94);
    \foreach \y in {89,85,81,77,73,69,65} {\draw[gray!70] (4,\y) -- (30,\y);}
    \node[font=\scriptsize\itshape] at (17,92) {verso};
    \draw[->,>=Stealth] (37,10) .. controls (46,30) and (44,52) .. (32,58);
    \node[font=\scriptsize,align=center] at (17,-6) {recto, puis verso\\du feuillet rabattu};
  \end{scope}
\end{tikzpicture}
\caption{Disposition matérielle du carnet relié en tête et rabattement des feuillets.}
\end{figure}

\subsection{Feuillets et foliotation}

Le carnet compte 93~feuillets écrits, soit 184~pages, plus le plat supérieur de la couverture. L'original n'est pas folioté : la foliotation de l'édition est restituée. Six pages sont écrites en travers, le carnet tourné d'un quart de tour (ff.~37~v\textsuperscript{o}--40~r\textsuperscript{o}, du 15 au 18~octobre 1916, pendant le mouvement de la compagnie, de Ribeaucourt vers Albert).

Le feuillet~6 (5~août 1916) n'est écrit qu'au recto : son verso, photographié, est blanc, et le récit reprend le 6~août au f.~7~r\textsuperscript{o}. Seul le recto du feuillet~93 (30~avril 1917) a été photographié : c'est la dernière page écrite connue, à la date annoncée par l'étiquette. La concordance des feuillets et des pages de l'édition figure en annexe.

\subsection{Disposition des pages}

Chaque journée commence, à gauche, par une date soulignée (jour de la semaine, souvent abrégé, et quantième) ; le mois et l'année sont rappelés, soulignés, en haut à droite. Le récit, très elliptique, est le plus souvent suivi des dépenses du jour, séparées par un trait horizontal et souvent réunies par un grand crochet : blanchissage (\textit{washing}), nourriture (\textit{food}), cartes postales, pourboires au palefrenier (\textit{tip au groom}), journaux, menus achats. Une page peut contenir plusieurs journées, séparées par un filet, et une journée occuper plusieurs pages. La transcription reproduit ces traits, crochets et accolades.

\subsection{Histoire du document et de la présente édition}

Sur la transmission du carnet et le projet des \textit{Feuillets de Route}, voir p.~\pageref{intro:transmission} et l'annexe, p.~\pageref{annexe:doc6}.

L'édition repose sur les recherches de Pierre Combes et de Pierre-Alain Hauseux. Elle donne tout le carnet et les pièces annexes : lecture de chaque page sur des photographies en haute définition, déchiffrement lettre à lettre des passages difficiles, restitution de la disposition et de l'ordre des feuillets, commentaire critique et historique. La composition et la révision philologique ont été faites avec l'aide d'outils informatiques et d'une super intelligence ; les lectures douteuses sont signalées.
